% SGA 6, Exposición XII: fragmento español alineado con la fuente.
% Alcance: sección 3 completa, «Démonstration de I : le dévissage de Oort».
% Fuente francesa: sga6_fr_r1/src/fr.tex, líneas 14324--14468.
% Las lecturas fuente manifiestamente defectuosas se conservan como primer
% argumento de \sisi; la lectura coherente usada por la edición es la segunda.

\providecommand{\sisi}[2]{#2}

\section*{3. Demostración de I: el dévissage de Oort}

\textbf{Proposición 3.1.} (Oort, Bull. soc. math. France, t. 90-1962).
Sean $S$ un esquema localmente noetheriano, $f:X\to S$ un morfismo propio y
$X'\subset X$ un subesquema cerrado. Se supone que:
\begin{enumerate}[label=\roman*)]
\item $X'$ está definido por un ideal $I$ contenido en el nilradical $N$ de
$X$ y $N.I=0$.
\item $X'$ es plano sobre $S$ y cohomológicamente plano sobre $S$ en
dimensión $0$ (EGA III 7.8.1).
\item $I$ es plano sobre $S$ y cohomológicamente plano sobre $S$ en
dimensiones $0$ y $1$, de modo que (según EGA III 7.4.6) existen un
$\mathcal O_S$-módulo coherente $Q$ y un isomorfismo funtorial en el
$\mathcal O_S$-módulo coherente $M$,
\[
R^2f_*(I\otimes f^*M)\simeq \Hom_{\mathcal O_S}(Q,M).
\]
\item El conúcleo $E$ de $f_*(\mathcal O_{X'})\to R^1f_*(I)$ es un
$\mathcal O_S$-módulo localmente libre.
\item El homomorfismo
$(f_*\mathcal O_X)_{\red}\longrightarrow(f_*\mathcal O_{X'})_{\red}$ es
sobreyectivo.
\end{enumerate}
Entonces el morfismo canónico
$\Pic_{X/S}\longrightarrow\Pic_{X'/S}$ es representable por un morfismo
afín de presentación finita; más precisamente, se tiene una sucesión exacta
canónica de haces f.p.p.f.,
\[
0\longrightarrow \mathbb V(E^\vee)\longrightarrow \Pic_{X/S}
\longrightarrow\Pic_{X'/S}\xrightarrow{u}\mathbb V(Q)
\]
y, si $P=\Ker(u)$, el monomorfismo $P\hookrightarrow\Pic_{X'/S}$ es
representable por una inmersión cerrada y $\Pic_{X/S}$ es un torsor sobre
$P$ bajo el grupo $\mathbb V(\check E)$ (es decir, la sucesión de haces
f.p.p.f.
\[
0\longrightarrow \mathbb V(\check E)\longrightarrow \Pic_{X/S}
\longrightarrow P\longrightarrow0
\]
es exacta).

Consideremos, para un $S$-esquema $T$, los dos diagramas siguientes, de
filas exactas:
\[
\begin{tikzcd}
0 \arrow[r] & I_T \arrow[r] & \mathcal O_{X_T} \arrow[r,dash] &
\mathcal O_{X'_T} \arrow[r] & 0\\
0 \arrow[r] & I_T \arrow[r,"\alpha",dash] \arrow[u,equal] &
\mathcal O_{X_T}^{*} \arrow[r] \arrow[u] &
\mathcal O_{X'_T}^{*} \arrow[r] \arrow[u] & 0 ,
\end{tikzcd}
\qquad (\alpha(i)=1+i)
\]
y
\[
\begin{tikzcd}[column sep=tiny]
(f_T)_*(\mathcal O_{X'_T}) \arrow[r,"\delta_T"] &
R^1(f_T)_*(I_T) \arrow[d,equal]\\
(f_T)_*(\mathcal O_{X'_T}^{*}) \arrow[r,"\delta_T^*"] \arrow[u] &
R^1(f_T)_*(I_T) \arrow[r] &
R^1(f_T)_*(\mathcal O_{X_T}^{*}) \arrow[r] &
\sisi{R'(f_T)_*}{R^1(f_T)_*}(\mathcal O_{X'_T}^{*}) \arrow[r] &
R^2(f_T)_*(I_T) .
\end{tikzcd}
\]
Adviértase que el cuadrado izquierdo de este diagrama \emph{no es
necesariamente conmutativo}. Sin embargo, se tiene:
\[
\textbf{3.2.}\qquad \im\delta_T^*=\im\delta_T.
\]
En efecto, se puede suponer evidentemente que $S$ y $T$ son afines. Tomemos
un recubrimiento afín $\{V_i\}$ de $X_T$. Sean
$g\in\Gamma(X_T;\mathcal O_{X'_T}^{*})$ y $g_i$ un levantamiento de
$g|V_i$ en $\Gamma(V_i;\mathcal O_{X_T}^{*})$. Entonces
$\delta_T^*(g)$ es igual a la clase del $1$-cociclo
$\delta_{ij}^*=(g_ig_j^{-1}-1)$; $\delta_T(g)$, a la de
$\delta_{ij}=(g_i-g_j)=g_j\delta_{ij}^*=g\,\delta_{ij}^*$, y por tanto
\[
\delta_T(g)=g\,\delta_T^*(g),
\]
teniendo en cuenta la acción de
$\Gamma(X_T,\mathcal O_{X'_T})$ sobre $I_T$ y, por consiguiente, sobre
$H^1(X_T,I_T)$.

Según ii), se tiene
$\Gamma(X_T;\mathcal O_{X'_T})=
\Gamma(X;\mathcal O_{X'})\otimes_S\Gamma(T,\mathcal O_T)$ y, por tanto,
$g=\sum g_\beta\otimes t_\beta$, con
$g_\beta\in\Gamma(X,\mathcal O_{X'})$ y
$t_\beta\in\Gamma(T,\mathcal O_T)$. Ahora bien, por (v),
$(f_*\mathcal O_X)_{\red}\longrightarrow
(f_*\mathcal O_{X'})_{\red}$ es sobreyectivo; por consiguiente,
$g_\beta=v(h_\beta)+n_\beta$, con
$h_\beta\in\Gamma(X;\mathcal O_X)$,
$v:\Gamma(X,\mathcal O_X)\to\Gamma(X,\mathcal O_{X'})$ y
$n_\beta\in\operatorname{nilrad}(\Gamma(X,\mathcal O_{X'}))$. Según i),
$(n_\beta\otimes t_\beta)I_T=0$. Por tanto,
$\delta_T(g)=h\,\delta_T^*(g)$, donde
$h=\sum v(h_\beta)\otimes t_\beta\in
\Gamma(X_T,\mathcal O_{X_T})$; de aquí
$\delta_T^*(g)=h^{-1}\delta_T(g)=\delta_T(h^{-1}g)$ y, por tanto,
$\im\delta_T^*\subset\im\delta_T$.

Finalmente, sea
$g=\sum g_\beta\otimes\sisi{t_\alpha}{t_\beta}
\in\Gamma(X_T,\mathcal O_{X'_T})$. Sea
$g_\beta=v(\sisi{h}{h_\beta})+n_\beta$, donde
$n_\beta\in\operatorname{nilrad}
(\Gamma(X,\sisi{\mathcal O_X}{\mathcal O_{X'}}))$; entonces
$\delta_T(g)=\delta_T(\sum n_\beta\otimes t_\beta)=
\delta_T(1+\sum n_\beta\otimes t_\beta)=
\delta_T^*(1+\sum n_\beta\otimes t_\beta)$, de donde
$\im\delta_T\subset\im\delta_T^*$.

En virtud de 3.2, iii) y iv), se tiene una sucesión exacta funtorial en $T$:
\[
0\longrightarrow \mathbb V(E^\vee)(T)\longrightarrow
R^1(f_\sisi{t}{T})_*(\mathcal O_{X_T}^{*})
\longrightarrow
R^1(f_T)_*(\sisi{\mathcal O_{X'_T}}{\mathcal O_{X'_T}^{*}})
\longrightarrow
\mathbb V(Q)(T),
\]
que, al pasar a los haces f.p.p.f. asociados, da una sucesión exacta
\[
0\longrightarrow \mathbb V(E^\vee)\longrightarrow\Pic_{X/S}
\longrightarrow\Pic_{X'/S}\xrightarrow{u}\mathbb V(Q).
\]
Si $P=\Ker(u)$, el monomorfismo $P\hookrightarrow\Pic_{X'/S}$ se obtiene
por cambio de base a partir de la sección unidad
$S\hookrightarrow\mathbb V(Q)$; como $\mathbb V(Q)$ es representable y
separado, resulta que $P\hookrightarrow\Pic_{X'/S}$ es representable por
una inmersión cerrada.

Finalmente,
$0\longrightarrow\mathbb V(E^\vee)\longrightarrow\Pic_{X/S}
\xrightarrow{w}P\longrightarrow0$ es un torsor para la topología f.p.p.f.;
a fortiori, gracias al descenso f.p.q.c. a partir de una «escisión» del
torsor, $w$ es representable por un morfismo afín fielmente plano de
presentación finita (véase Oort, «~Commutative group schemes~»,
Proposición (17.4)).

\textbf{Corolario 3.3.} Bajo las condiciones de 3.1, se tiene:
\begin{enumerate}[label=\roman*)]
\item Si
$\Pic_{(\zar)(X'/S)}\to\Pic_{(\fppf)(X'/S)}$ es un isomorfismo, entonces
$\Pic_{(\zar)(X/S)}\to\Pic_{(\fppf)(X/S)}$ es un isomorfismo.
\item Suponiendo que el $S$-esquema $T$ es afín, si
$\Pic(X'_T)\to\Pic_{(\zar)(X'/S)}^{(T)}$ es sobreyectivo, entonces
$\Pic(X_T)\to\Pic_{(\zar)(X/S)}^{(T)}$ es sobreyectivo.
\end{enumerate}
En efecto, (i) resulta del diagrama siguiente de haces para la topología de
Zariski,
\[
\begin{tikzcd}[column sep=small]
0 \arrow[r] \arrow[d,"S"] & \mathbb V(E^\vee) \arrow[r] \arrow[d,"S"] &
\Pic_{(\zar)(X/S)} \arrow[r] \arrow[d] &
\Pic_{(\zar)(X'/S)} \arrow[r] \arrow[d,"S"] &
\mathbb V(Q) \arrow[d,"S"]\\
0 \arrow[r] & \mathbb V(E^\vee) \arrow[r,dash] &
\Pic_{(\fppf)(X/S)} \arrow[r] &
\Pic_{(\fppf)(X'/S)} \arrow[r] &
\mathbb V(Q),
\end{tikzcd}
\]
aplicándole el lema de los cinco.

(ii) resulta del siguiente diagrama de filas exactas:
\[
\begin{tikzcd}[column sep=small]
H^1(X_T,I_T) \arrow[r] \arrow[d] &
H^1(X_T,\mathcal O_{X_T}^{*}) \arrow[r,dash] \arrow[d] &
H^1(X_T,\mathcal O_{X'_T}^{*}) \arrow[r,dash] \arrow[d] &
H^2(X_T,I_T) \arrow[d]\\
H^0(T,R^1(f_T)_*(I_T)) \arrow[r] &
H^0(T,R^1(f_T)_*\mathcal O_{X_T}^{*}) \arrow[r] &
H^0(T,R^1(f_T)_*\mathcal O_{X'_T}^{*}) \arrow[r] &
H^0(T,R^2f_*I_T),
\end{tikzcd}
\]
pues las flechas verticales exteriores son isomorfismos, ya que $T$ es
afín.

\textbf{Corolario 3.4.} Sean $k$ un cuerpo y $X$ un $k$-esquema propio. Se
supone que $\Pic_{X_{\red}/k}$ posee un haz invertible universal. (Por
ejemplo, este será el caso, según la teoría de descenso, si toda componente
conexa $X'_i$ de $X_{\red}$ posee un punto racional; por ejemplo, si $k$ es
algebraicamente cerrado o, más generalmente, si el m.c.d. de los grados de
los $0$-ciclos de todo $X'_i$ es $1$). Entonces, sobre todo abierto afín
$V$ de $\Pic_{X/k}$, existe un haz invertible universal sobre
$X\times_kV$.

Finalmente, se demuestra I en la forma más precisa siguiente:

\textbf{Proposición 3.5.} Sean $S$ un esquema noetheriano íntegro y
$f:X\hookrightarrow Y$ una inmersión nilpotente de $S$-esquemas propios.
Entonces existe un abierto no vacío $V$ de $S$ tal que el morfismo
\[
f^*:\Pic_{Y/S}|V\longrightarrow\Pic_{X/S}|V
\]
sea representable por un morfismo afín de presentación finita; de hecho,
$N=\Ker f^*$ es un esquema de grupos afines con fibras unipotentes (en el
sentido de SGA 3 XVII).

En efecto, $X$ está definido por un ideal nilpotente $I$. Entonces $I$
posee una filtración por ideales
$I_0=I\supset I_1\supset\cdots\supset I_n=0$ tales que
$NI_i\subset I_{i-1}$, donde $N$ es el nilradical de $\mathcal O_Y$ (por
ejemplo, $I_k=N^kI$). Sea $X_i$ el subesquema de $Y$ definido por $I_i$.
Por recurrencia, basta estudiar las inmersiones
$X_i\hookrightarrow X_{i-1}$. Se puede suponer, por tanto, que $N.I=0$;
esta es la condición i) de 3.1, y se concluye reduciéndose al caso en que se
verifican también las demás condiciones (ii) a (v). Gracias al teorema de
planitud genérica (EGA IV 6.9.1), aplicado\sisi{}{, con las notaciones de 3.1,} a $\mathcal O_{X'}$, a los % Canon SGA-SRC-001315
$R^if_*(I)$ y a $E$, y teniendo en cuenta las relaciones de Künneth EGA III
6.9.9, se puede suponer, tras restringir $S$, que se verifican las
condiciones ii), iii) y iv).

Existe una extensión finita $K$ de $k(S)$ tal que
$(f_K)_*(\mathcal O_{X'_K})$ tenga todas sus extensiones residuales
triviales. Si $S'$ es un esquema íntegro finito y plano sobre un abierto
$V$ de $S$, con cuerpo de funciones $k(S')=K$, existe un abierto $V'$ de
$S'$ tal que
\[
\bigl[(f_{S'})_*(\mathcal O_{X'_{S'}})\bigr]_{\red}
\simeq
\bigl[(f_{S'})_*(\mathcal O_{X_{S'}})\bigr]_{\red}
\]
sobre $V'$; tras restringir $S$ y $S'$, se puede suponer $S=V$ y $S'=V'$.
Como la afirmación que se ha de demostrar es local para la topología
f.p.q.c., se puede reemplazar $S$ por $S'$. Finalmente, se aplica 3.1.

\textbf{3.6. Observación.} Se deduce asimismo, teniendo en cuenta las
sucesiones exactas
$\Pic_{X_{i-1}/S}\longrightarrow\Pic_{X_i/S}\longrightarrow\mathbb V(Q_i)$,
que para todo $s\in S$
\[
\coker\bigl(\Pic^0_{X_s/s}\longrightarrow\Pic^0_{Y_s/s}\bigr)
\]
es un grupo unipotente.
